\documentclass[11pt]{article}
\usepackage[margin=1in]{geometry}
\usepackage{booktabs}
\usepackage{hyperref}

\title{CAI6734 Weekly Report 8: Locating the Bottleneck and Separating the Two Agents}
\author{Omer Kahveci}
\date{November 17, 2026}

\begin{document}
\maketitle

\section{Project Objective}

Week 7 ended with a full-pass baseline that matched the paper (85.4\% test
accuracy at \$1,259 per case) and two open weaknesses: the Decision Agent (DA)
drifted toward ordering no tests, and the Hypothesis Agent's (HA's) confidence
did not improve during training. I planned to review the literature on agentic
search and to pick one improvement to the agentic loop.

This week I turned that plan into a sequence of three hypotheses, each
pre-registered in writing (prediction, primary comparison and decision rule)
before any result existed:

\begin{itemize}
  \item \textbf{H1:} does the DA's \emph{choice} of tests matter, or only how
        many it orders?
  \item \textbf{H2:} is an HA trained on its own, then frozen, a better
        diagnoser than the HA inside the shared LA-CDM adapter?
  \item \textbf{H3:} does a DA trained against that frozen HA perform better,
        with the original reward (H3a) and with a reward for informative tests
        (H3b)?
\end{itemize}

\section{Accomplishments}

\subsection*{My Contributions This Week}

\begin{itemize}
  \item Wrote a literature review on reinforcement learning for multi-turn,
        information-seeking agents and sequential diagnosis, which motivated
        H1--H3.
  \item Finished the third baseline run (attempt 3, with last week's HA
        training fix) and tested it.
  \item Built scripted test-ordering policies (no tests, random tests, all
        tests) for H1 and ran them on two trained adapters.
  \item Built a standalone HA training and scoring pipeline for H2, trained two
        seeds, and scored five HAs on identical patient records.
  \item Integrated the frozen HA into the LA-CDM environment as a second LoRA
        adapter on the same model, with unit tests checking that the DA's
        transcript format is unchanged.
  \item Implemented DA-only training against the frozen HA and an
        information-gain reward for H3, and ran both arms (900 steps each).
\end{itemize}

All comparisons below use the same 240 test patients, paired tests (exact
McNemar for accuracy, paired bootstrap for intervals) and 95\% confidence
intervals, as in Week 7. Checkpoints were always chosen on validation.

\textbf{Names used below.}
\begin{itemize}
  \item \textbf{Untrained model:} Llama-3-8B-Instruct with no training on this
        task.
  \item \textbf{Baseline model:} our LA-CDM reproduction from Week 7 (85.4\%
        accuracy). Its two roles share one LoRA adapter: the \emph{baseline DA}
        and the \emph{baseline HA}.
  \item \textbf{New HA:} the HA trained on its own in H2 (seed 0 unless
        stated), then frozen.
  \item \textbf{Retrained DA:} a new DA trained against the frozen new HA in
        H3, with the original reward (H3a) or with an added belief reward (H3b).
\end{itemize}

\subsection*{Attempt 3: The HA Training Fix Alone}

Attempt 3 repeated the baseline model's training with only the HA loss fix. Its step-600
checkpoint, chosen on validation, reached 75.4\% test accuracy at \$1,115 per
case, about 10 points below the baseline model. Its HA almost never diagnosed
pancreatitis correctly (44\% recall). With one run per code version, this gap
cannot be attributed to the fix; it is at least partly run-to-run variance. That
variance is itself the main finding: two nearly identical training runs gave
HAs that differ by 7 points (see H1 below).

\subsection*{H1: The DA's Choice of Tests Matters}

The literature suggested that a random test planner paired with a strong
diagnoser can beat LA-CDM on this benchmark (LDTL, 2026). H1 tested this on our
model. I kept the trained HA and replaced only the DA's test choices, scoring
every arm by the HA's final hypothesis. The rule, fixed in advance, was that H1
would be supported if the interval for ``DA tests minus random tests'' lay
inside $\pm$5 points.

\begin{table}[ht]
  \centering
  \caption{H1: accuracy of the baseline HA's last hypothesis given different test
  choices. Only the tests change between rows; every row is scored by the HA, not by
  the DA. The first row is therefore 87.1\%, not the baseline model's 85.4\%, which
  is the DA's final diagnosis.}
  \label{tab:h1}
  \small
  \begin{tabular}{lrrr}
    \toprule
    Tests given to the HA & HA accuracy & Tests per case & Cost per case \\
    \midrule
    Chosen by the baseline DA & 87.1\% & 1.19 & \$1,259 \\
    None (history only) & 62.9\% & 0 & \$0 \\
    Random, same number per patient & 64.6\% & 1.19 & \$494 \\
    All available tests & 86.2\% & 9.30 & \$4,082 \\
    \bottomrule
  \end{tabular}
\end{table}

\textbf{H1 was rejected.} The DA's tests beat random tests at the same count by
22.5 points [+16.7, +28.3] ($p < 0.001$), and attempt 3 replicated the result
(+9.6 [+4.6, +14.6]). The DA also matched ordering every test at about a third
of the cost. However, the ``all tests'' arm revealed the real limit: even with
every result, the baseline HA reached only 86.2\% and attempt 3's HA 79.2\%,
against about 94\% for LDTL's separately trained diagnoser. Better search cannot
raise accuracy above the HA's ceiling.

\textbf{A deviation, stated openly.} The H1 rule said to build the search reward
next. Because the HA's ceiling caps any search improvement, I did H2 (a better
HA) first. I changed the plan after seeing the H1 results and recorded it as
such; the rule itself was not changed.

\subsection*{H2: A Standalone HA Is a Better Diagnoser}

In LA-CDM both agents share one LoRA adapter and are trained in turns, so
updates for one role can undo the other's, and the HA learns only from records
the DA happens to produce. For H2 I trained a separate LoRA adapter for the HA
alone, with supervised fine-tuning on four views of every training patient:
history only, all tests, and two random subsets of tests. The adapter has the
same size as the shared one, so a gain is not due to more parameters.

Generating answers made HAs hard to compare, so I scored every HA the same way:
its answer is the disease with the highest probability, and its confidence is
that probability. This gives a full probability distribution for every case.
Each HA read exactly the same test results, copied from the H1 runs.

\begin{table}[ht]
  \centering
  \caption{H2: accuracy of each HA on identical records. Columns give the test
  results shown to every HA for each patient: none; random tests (as many as the
  baseline DA ordered); the tests the baseline DA ordered; or every available
  test. Every row is scored by its highest-probability disease. ECE is with all
  tests. The main comparison is the baseline HA against the new HA (seed 0).}
  \label{tab:h2}
  \small
  \begin{tabular}{lrrrrr}
    \toprule
    HA & No tests & Random tests & Baseline DA's tests & All tests & ECE \\
    \midrule
    Untrained model (zero-shot) & 61.3\% & 63.3\% & 82.5\% & 75.4\% & 10.8\% \\
    Baseline HA (shared adapter) & 63.3\% & 64.6\% & 87.9\% & 87.9\% & 5.6\% \\
    New HA, seed 0 & 78.3\% & 80.8\% & 88.8\% & 90.0\% & 4.5\% \\
    New HA, seed 1 & 76.2\% & 81.7\% & 91.7\% & 93.8\% & 1.2\% \\
    \bottomrule
  \end{tabular}
\end{table}

Three notes on Table~\ref{tab:h2}:
\begin{itemize}
  \item The baseline HA scores 87.9\% on the baseline DA's tests here. That is
        higher than the baseline model's 85.4\%, which is the DA's final diagnosis,
        and its generated hypothesis was right 87.1\% of the time (Table~\ref{tab:h1}).
        The differences come from the DA overriding the HA and from reading
        probabilities instead of generated text.
  \item The untrained model does better with the baseline DA's tests than with
        all tests. This is likely because the all-tests records are long (about
        nine tests, mostly lab panels), and an untrained model reads one
        focused imaging report more reliably. I have not verified this. Every
        trained HA does as well or better with all tests.
  \item The zero-shot row is scored by probabilities, so it cannot fail the
        answer format. It is not comparable to the Week 7 zero-shot system
        (8.8\%), where the untrained model also acted as the DA and mostly
        failed the format.
\end{itemize}

\textbf{H2 was supported for partial records.} The new HA (seed 0) beat the baseline
HA by 15.0 points [+8.7, +21.3] with no tests and 16.2 points [+10.0, +22.1]
with random tests. With all tests the gain was +2.1 [$-$2.1, +6.2], not
detectable. It was also better calibrated and no longer over-called
pancreatitis. One pre-registered check failed: the two seeds differ by 3.7
points with all tests ($p = 0.022$), although their validation losses were
almost identical, so validation could not have told them apart. Seed 0 was
fixed in advance as the HA to freeze, and it is used in everything below. Finally,
the new HA differs from the baseline HA in both its adapter and its training data, so
this result does not separate the effect of the separate adapter from that of
training on partial records.

\subsection*{Combining the Baseline DA with the Frozen New HA}

Before retraining anything, I placed the frozen new HA inside the environment
next to the baseline DA, as a second LoRA adapter on the same model. The DA's
transcript has exactly the same format as before; only the source of the
hypothesis changes.

\begin{table}[ht]
  \centering
  \caption{The baseline DA with the baseline HA and with the frozen new HA.}
  \label{tab:step2}
  \small
  \begin{tabular}{lrrr}
    \toprule
    Metric & Baseline model & Baseline DA + new HA & Difference [95\% CI] \\
    \midrule
    Accuracy, DA's diagnosis & 85.4\% & 85.4\% & +0.0 [$-$3.8, +3.7] \\
    Accuracy, HA's final hypothesis & 87.1\% & 89.2\% & +2.1 [$-$1.7, +5.8] \\
    Cost per case & \$1,259 & \$1,285 & +\$27 [$-$\$30, +\$91] \\
    Calibration error (ECE) & 14.8\% & 5.2\% & \\
    AUROC, confidence vs.\ correct & 0.750 & 0.936 & \\
    \bottomrule
  \end{tabular}
\end{table}

Accuracy and cost did not change, but the confidence the DA reads became far
more informative. The key observation is that the HA's final hypothesis was
right 89.2\% of the time while the DA's diagnosis was right 85.4\% of the time:
the DA overrides correct answers, mostly on pancreatitis (HA recall 83.3\%,
DA recall 68.5\%). The DA had been trained alongside the baseline HA, whose confidence was
inflated, which is a plausible reason it does not rely on the hypothesis.

\subsection*{H3: Retraining the DA Against the Frozen HA}

H3 trains a new DA from the untrained model while the new HA stays frozen. Two arms
separate the effect of the partner from the effect of the reward:

\begin{itemize}
  \item \textbf{H3a:} the original reward (diagnosis, format and test cost).
  \item \textbf{H3b:} the same reward plus the frozen HA's final probability of
        the correct disease. GRPO compares rollouts of the same patient, and
        the HA's belief before any test is the same for all of them, so this
        term trains exactly like the total information gained by the tests. A
        test therefore pays off when it raises that probability by more than
        its cost.
\end{itemize}

The decision rule was fixed before any H3 result: an arm is supported only if the
95\% interval of its paired accuracy gain is entirely above zero. H3a is compared
with the baseline DA + new HA (the same HA, but a DA trained with the baseline HA).
H3b is compared with the baseline DA + new HA and with H3a.

Both arms trained for 900 steps. By validation accuracy, step 700 was the best
checkpoint in both (73.8\% for H3a and 70.4\% for H3b), and I tested only those
two checkpoints. Table~\ref{tab:h3} gives the test results.

\begin{table}[ht]
  \centering
  \caption{H3 test results (240 cases). All runs except the baseline model use
  the frozen new HA. Accuracy is shown for the DA's final diagnosis and for the
  HA's last hypothesis.}
  \label{tab:h3}
  \small
  \begin{tabular}{lrrrr}
    \toprule
    Run & DA's diagnosis & HA's last hypothesis & Tests per case & Cost per case \\
    \midrule
    Baseline model & 85.4\% & 87.1\% & 1.19 & \$1,259 \\
    Baseline DA + new HA & 85.4\% & 89.2\% & 1.20 & \$1,285 \\
    Retrained DA (H3a) & 85.8\% & 88.3\% & 1.27 & \$1,025 \\
    Retrained DA + belief reward (H3b) & 74.6\% & 78.3\% & 2.10 & \$536 \\
    \bottomrule
  \end{tabular}
\end{table}

\textbf{H3a: no detectable accuracy gain, but cheaper.} Compared with the baseline DA + new HA,
accuracy changed by +0.4 points [$-$2.5, +3.3] ($p = 1.00$), so under the
pre-registered rule H3a is not an improvement. Cost per case fell by \$260
[$-$\$348, $-$\$182] ($p < 0.001$), about 20\%, a secondary outcome. H3a kept
ordering CT (66.7\% of patients against 68.3\% for the baseline DA) but nearly stopped
ordering ultrasound (3.8\% against 20.0\%), using blood panels instead. The gap
between the DA's diagnosis and the HA's final hypothesis narrowed from 3.8 to
2.5 points, which suggests the retrained DA relies slightly more on the HA.

\textbf{H3b: worse.} The belief reward lowered accuracy by 10.8 points
[$-$16.7, $-$5.4] against the baseline DA + new HA and by 11.2 points [$-$16.2, $-$6.2] against
H3a (both $p < 0.001$), while cutting cost to \$536 per case. Test rates show
why: H3b almost stopped ordering CT (1.7\% of patients against 66.7\% for H3a)
and ordered cheap tests instead, such as a complete blood count (74.2\%), an
electrolyte panel (35.8\%) and a radiograph (33.3\%). The HA's own final
hypothesis fell to 78.3\%, so the loss comes from less informative tests, not
from the DA ignoring the HA. My interpretation, which I have not verified, is
that the frozen HA is often fairly confident after the history alone. A CT then
raises its probability of the correct disease by less than the CT's cost, while
several cheap tests each move it slightly for almost no cost. The reward
measures marginal changes in the HA's belief, which is not the same as getting
the diagnosis right.

\subsection*{What I Learned}

Asking narrow questions in order located the bottleneck. H1 showed the search
was already good, and that the diagnoser set the ceiling. H2 showed that an HA
trained on its own, on partial records, diagnoses far better from thin records.
Combining it with the baseline DA showed that a better
diagnoser alone is not enough when the DA was trained not to trust it, and H3a
showed that retraining the DA against it lowered cost rather than raising
accuracy. H3b showed that a reward which sounds aligned with good diagnosis,
information gained by the diagnoser, can be satisfied more cheaply by many
uninformative tests than by one decisive test. Writing
each decision rule in advance kept me from reading the results the way I hoped,
for example when H2's stability check failed. Two near-identical runs can
differ by 7 to 10 points, so single-run comparisons between code versions
(the baseline model against attempt 3) need to be treated with care.

\section{Challenges}

DA-only training was unstable late in both H3 runs. In H3a the DA nearly
stopped ordering tests at step 800 (0.34 per case), and at step 900, 95\% of
validation cases were undiagnosed, so the automatic collapse guard stopped the
run. H3b dropped to 54\% undiagnosed at step 800 and then recovered. I checked
the job logs and ruled out the restart after the 24-hour time limit as the
cause. Both step-700 checkpoints come before the instability, but it limits how
long the DA can be trained at this learning rate.

H3a's cost reduction is a secondary outcome; the pre-registered primary outcome,
accuracy, did not change. All H2 and H3 results use one HA seed and one DA seed. The two HA seeds differ
significantly, so the H3 results may depend on which HA was frozen. As before,
the 240-patient test set gives intervals of about $\pm$4--6 points for paired
differences, and the diverticulitis class has only 25 test cases.

\section{Next Steps}

\begin{itemize}
  \item Inspect H3b transcripts to confirm why CT stopped being ordered, and
        decide whether a reweighted belief reward (a smaller weight, or a larger
        reward for a correct diagnosis) is worth one more run. Per-test credit
        (as in IGPO) would likely face the same problem, because it uses the same
        measure of gain.
  \item Investigate the late-training instability, for example with a lower
        learning rate or a stronger KL penalty, before any longer runs.
  \item Repeat H3a with the second HA seed to check whether its cost reduction
        holds, since the two HA seeds differ significantly.
  \item Start assembling the final report around the main finding so far: the
        shared adapter limits the HA, and separating the two agents gives a
        better-calibrated system at equal accuracy and lower cost.
\end{itemize}

\section{AI Disclaimer}

I used Anthropic Claude to assist with designing the experiments, implementing
the code, inspecting training and evaluation logs, computing the statistics,
and drafting this report.

\end{document}
